\documentclass{article}
\usepackage{amsmath}
\usepackage{amssymb}
\usepackage{enumitem}
\usepackage{algorithm}
\usepackage{listings}
\usepackage{color,xcolor}
\usepackage[T1]{fontenc}
\usepackage{etoolbox}
\usepackage{multicol}
\usepackage{geometry}
\usepackage[colorlinks=true,linkcolor=blue,urlcolor=red,bookmarksopen=true]{hyperref}
\usepackage{tikz, pgfplots, tkz-euclide,calc}
    \usetikzlibrary{patterns,snakes,shapes.arrows,3d,patterns.meta,angles,quotes}
    \geometry{
        total = {160mm, 237mm},
        left = 25mm,
        right = 35mm,
        top = 30mm,
        bottom = 30mm,
      }

\usepackage{tcolorbox}
     \tcbuselibrary{listings,skins}

\newcommand{\enter}{\raisebox{-1.8pt}{\begin{tikzpicture}[scale=0.3]
    \draw[thin,fill=lightgray] (0,0) rectangle (2,1);
    \draw (0.3,0.3) -- (0.7,0.3)--(0.7,0.6);     
\end{tikzpicture}}}

\definecolor{HIMAmuda}{HTML}{01D1FD}
\definecolor{HIMAtua}{HTML}{02016A}
\definecolor{HIMAabu}{HTML}{CBCBCC}
\definecolor{pgray}{rgb}{0.5,0.5,0.5}
\definecolor{pblue}{rgb}{0.13,0.13,1}
\definecolor{pgreen}{rgb}{0,0.5,0}
\definecolor{pred}{rgb}{0.9,0,0}
\definecolor{pgrey}{rgb}{0.46,0.45,0.48}
\definecolor{pcyan}{HTML}{D4EFFC}
\definecolor{lblue}{HTML}{00AEEF}
\definecolor{input}{HTML}{AAE1FA}
\definecolor{bg}{rgb}{0.95, 0.95, 0.92}
\definecolor{vscode}{HTML}{282A36}
\definecolor{PastelGreen}{HTML}{77DD77}

\newcommand{\inputscan}[1]{\raisebox{0pt}[1pt]{\colorbox{darkgray}{#1}}}

\usepackage{listings}

\lstdefinestyle{Liang}{
language=Java,
showspaces=false,
showtabs=false,
breaklines=true,
showstringspaces=false,
breakatwhitespace=true,
commentstyle=\color{pgray},
keywordstyle=\color{pblue},
stringstyle=\color{pgreen},
basicstyle=\small\ttfamily,
frame=single,
backgroundcolor=\color{pcyan},
escapeinside={(*}{*)},}

\lstdefinestyle{output}{
    language=Java,
    backgroundcolor=\color{vscode},
    basicstyle=\small\ttfamily\color{white},
    frame=none,
    escapeinside={(*}{*)},
    showspaces=false,
    showtabs=false,
    breaklines=true,
    showstringspaces=false,
    breakatwhitespace=true,
    keywordstyle=\color{white},
    }

\lstdefinestyle{standard}{
    language=Java,
    showspaces=false,
    showtabs=false,
    breaklines=true,
    showstringspaces=false,
    breakatwhitespace=true,
    commentstyle=\color{pgray},
    keywordstyle=\color{pblue},
    stringstyle=\color{pgreen},
    basicstyle=\small\ttfamily,
    frame=single,
    backgroundcolor=\color{bg},
    escapeinside={(*}{*)},}
\lstset{style=Liang}

\newtcblisting{RunCode}[1][enhanced,drop shadow]{
    arc=0pt, outer arc=0pt,
    boxsep=1pt,
    boxrule=2pt,
    auto outer arc,
    colback=vscode,
    colframe=bg,
    listing only, 
    listing style=output,
    title=\color{black}Ex. Output,
    #1
    }

\newtcolorbox{hint}[1][]{
    colback=PastelGreen!5!white, 
    colframe=PastelGreen!75!black,
    fonttitle=\bfseries, 
    colbacktitle=PastelGreen!85!black,
    enhanced, 
    attach boxed title to top left={yshift=-2mm}, 
    title=Hint,
    #1
}

\newtcolorbox{req}[1][]{
    colback=lblue!5!white, 
    colframe=lblue!75!black,
    fonttitle=\bfseries, 
    colbacktitle=lblue!85!black,
    enhanced, 
    attach boxed title to top left={yshift=-2mm}, 
    title=Input,
    #1
}

\newtcolorbox{out}[1][]{
    colback=HIMAtua!5!white, 
    colframe=HIMAtua!75!black,
    fonttitle=\bfseries, 
    colbacktitle=HIMAtua!85!black,
    enhanced, 
    attach boxed title to top left={yshift=-2mm}, 
    title=Output,
    before upper=\renewcommand\thempfootnote{\Roman{mpfootnote}},
    #1
}

\renewcommand{\thesubsection}{\arabic{subsection}}
\newcommand{\R}{\mathbb{R}}
\newcommand{\Z}{\mathbb{Z}}

\title{\textbf{Tugas Pertemuan 4}}
\date{2 Oktober 2024}
\author{Dhanar A \& Fajar A}
\begin{document}
\maketitle

\section*{Problems}
\begin{enumerate}
    \pagenumbering{gobble}
    
    \item Buatlah sebuah program bilangan prima dengan batasan dari input user (Gunakan tipe boolean untuk pemberhentian program ini) 

    \begin{hint}
\begin{lstlisting}[style=standard]
final int perbaris = 10;
int count = 0;
int number = 2;
    \end{lstlisting}
    \end{hint}

\begin{RunCode}
Bilangan prima dalam 50 adalah 
2 3 5 7 11 13 17 19 23 29
31 37 41 43 47 53 59 61 67 71
73 79 83 89 97 101 103 107 109 113
127 131 137 139 149 151 157 163 167 173
179 181 191 193 197 199 211 223 227 229
\end{RunCode}
    \newpage
    
    \item Buatlah program countdown dengan catatan dd : hh : mm : ss 
    \begin{hint}
        \begin{itemize}
            \item gunakan \texttt{Thread.sleep(1000)} untuk menunggu 1 detik
            \item 1 hari = $\dfrac{\text{Target waktu}}{86400}$
            \item 1 jam = $\dfrac{(\text{Target waktu} \quad \texttt{\%} \quad 86400)}{3600} $
        \end{itemize}
    \end{hint}

    \begin{RunCode}
00:00:00:05
00:00:00:04
00:00:00:03
00:00:00:02
00:00:00:01
Waktu Habis !
    \end{RunCode}

    \newpage
    
    \item Sebuah perusahaan ingin menghitung bonus akhir tahun yang akan diberikan kepada karyawan berdasarkan jabatan mereka. Jabatan dan persentase bonus yang diberikan adalah sebagai berikut : 
    \begin{itemize}
        \item Manager : $20\%$ dari gaji pokok
        \item Supervisor : $15\%$ dari gaji pokok
        \item Staff : $10\%$ dari gaji pokok
        \item Intern : $5\%$ dari gaji pokok
    \end{itemize}

    Buatlah program yang menerima input berupa jabatan dan gaji pokok, kemudian menghitung bonus berdasarkan aturan diatas. Program juga harus memberikan pesan jika jabatan yang dimasukkan tidak valid .

    \begin{req}
        \begin{itemize}
            \item $0\leq \text{Gaji Pokok} \leq 100,\quad x_0,y_0\in\Z$ (Dalam Juta)
        \end{itemize}
    \end{req}
    \begin{out}
       \begin{itemize}
                \item $B:=$ Bonus untuk tiap jabatan\footnote{Cukup tampilkan 1 angka di belakang koma}\\
                $B\geq 0,\quad B\in\R$
            \end{itemize}      
    \end{out}
    \begin{hint}
        \begin{itemize}
            \item Gunakan \texttt{toLowerCase()} agar input mengabaikan huruf besar atau kecil 
            \item Tambahkan \texttt{Default} untuk jabatan yang tidak valid 
        \end{itemize}
    \end{hint}

    \begin{RunCode}
Masukkan gaji pokok : (*\inputscan{52} \enter*)
Masukkan jabatan : (*\inputscan{Manager} \enter*)
Gaji setelah ditambahkan bonus adalah 62.4 Juta
    \end{RunCode}

    
\end{enumerate}


\end{document}
